% ============================================================
% §9  Emergencja pól cechowania z substratu
% ============================================================
\section{Emergencja p\'ol cechowania}\label{sec:cechowanie}

Dotychczasowe sekcje wyprowadziły skalarny sektor TGP: pole $\Phi$
i~jego dynamikę. Niniejsza sekcja adresuje pytanie, jak
\textbf{pola cechowania} (elektromagnetyzm, oddziaływanie słabe
i~silne) wynikają z~tej samej ontologii --- emergencji z~substratu ---
nie zaś z~zewnętrznego postulatu. Kluczowa zasada:

\begin{remark}[Zasada emergencji cechowania]\label{rem:gauge-emergence-principle}
W~TGP pole $\Phi$ wynika z~\emph{modułu} parametru porządku
substratu. Pola cechowania wynikają z~\emph{fazy} tego parametru.
Obie wielkości są aspektami tej samej konfiguracji substratowej ---
nie są osobnymi bytami.
\end{remark}

Podkreślamy ontologiczną różnicę od standardowego podejścia:
nie wstawiamy pól cechowania na gotowej rozmaitości.
Wykazujemy, że w~granicy ciągłej substratu o~lokalnej symetrii
grupowej emergują \emph{spójne} pola wektorowe identyczne
ze standardowymi polami cechowania.

\begin{remark}[Status niniejszej sekcji]\label{rem:gauge-section-status}
Niniejsza sekcja ma status \textbf{otwartego programu badań},
nie zamkniętego modułu teorii. Precyzując poziomy:
\begin{itemize}[nosep]
  \item \textbf{Sformalizowany szkic}: emergencja U(1) z~fazy substratu
    (\S\ref{ssec:u1-photon}); kwantyzacja ładunku z~topologii wirów
    (stw.~\ref{prop:charge-quantization-gauge});
    niezmienniczość cechowania (krok~4 w~\S\ref{ssec:u1-photon}).
  \item \textbf{Twierdzenie topologiczne}: SU(2)$\times$U(1) i~SU(3)
    z~hierarchii defekt\'ow (\S\ref{ssec:su2-su3};
    formalny dow\'od: dod.~\ref{app:D2-defect-hierarchy},
    tw.~\ref{thm:D2-hierarchy-main}).
  \item \textbf{Program do domknięcia}: obliczone stałe sprzężenia
    $g_s, g_w, g'$ z~parametrów substratu; masy cząstek ze struktury
    defektów; chiralność fermionów.
\end{itemize}
Sekcja ta \textbf{nie powiela bytów}: pola cechowania są projekcją
\emph{fazowego} sektora tego samego substratu, który daje $\Phi$
w~sektorze amplitudowym (uwaga~\ref{rem:jeden-substrat} w~\S\ref{sec:ontologia}).
\end{remark}

% ----------------------------------------------------------
\subsection{Pe\l{}na struktura substratu: amplituda i~faza}%
\label{ssec:complex-substrate}

Sekcje~\ref{sec:pole}--\ref{sec:formalizm} wyprowadzi\l{}y
pe\l{}ny sektor grawitacyjny z~substratu relacyjnego $\Gamma = (V,E)$
z~realnym operatorem $\hat{s}_i \in \mathbb{R}$ i~symetri\k{a} $\mathbb{Z}_2$.
Ten substrat jest \textbf{kompletny} dla opisu grawitacji,
kosmologii i~fal grawitacyjnych.

Niniejsza sekcja rozszerza model o~\textbf{dodatkowy stopie\'n swobody}:
lokaln\k{a} faz\k{e} $\theta_i$, kt\'ora czyni substrat zespolonym.
Jest to \textbf{oddzielny krok badawczy} (por.\ uwaga~\ref{rem:jeden-substrat}
w~\S\ref{sec:ontologia}): rozszerzenie sektora grawitacyjnego,
nie zmiana jego fundamentu.
Pytanie: \emph{je\.zeli} substrat mo\.ze posiada\'c lokaln\k{a} faz\k{e},
czy emerguj\k{a} z~niej pola cechowania?

\begin{axiom}[Rozszerzenie substratu o~lokaln\k{a} faz\k{e}]%
\label{ax:complex-substrate}
Rozszerzamy substrat o~dodatkowy stopie\'n swobody:
ka\.zdy w\k{e}ze\l{} $i \in V$ nosi stan zespolony
(zamiast $\hat{s}_i \in \mathbb{R}$):
\begin{equation}\label{eq:complex-state}
    \psi_i \in \mathbb{C},
    \qquad
    \psi_i = \phi_i\, e^{i\theta_i},
    \qquad
    \phi_i = |\psi_i| \geq 0,
\end{equation}
gdzie $\phi_i$ jest \textbf{amplitudą} (intensywnością obecności),
a~$\theta_i \in [0, 2\pi)$ --- \textbf{fazą lokalną}.
\end{axiom}

Hamiltoniano pełnego substratu fazowego:
\begin{equation}\label{eq:H-complex}
    H_\Gamma^{(U1)} = \sum_i \left[
        \frac{m_0^2}{2}\,|\psi_i|^2
        + \frac{\lambda_0}{4}\,|\psi_i|^4
    \right]
    - J\sum_{\langle ij\rangle} \mathrm{Re}\!\left(
        \psi_i^*\,\psi_j
    \right).
\end{equation}
Ten hamiltoniano jest niezmienniczy względem globalnej transformacji
$U(1)$: $\psi_i \to e^{i\alpha}\psi_i$ ($\alpha = \mathrm{const}$).

\begin{remark}[Rozszerzenie jest minimalne i~odwracalne]\label{rem:amplitude-sector}
Dla $\theta_i \equiv 0$ hamiltoniano~\eqref{eq:H-complex}
redukuje si\k{e} dok\l{}adnie do~\eqref{eq:B-H} z~identyfikacj\k{a}
$\hat{s}_i = \phi_i$.
Sekcje~\ref{sec:pole}--\ref{sec:formalizm} s\k{a} zatem
\textbf{pe\l{}nym} i~\textbf{zamkni\k{e}tym} sektorem grawitacyjnym
--- nie wymagaj\k{a} fazy $\theta_i$.
Rozszerzenie o~$\theta_i$ jest motywowane pytaniem
o~pola cechowania i~jest logicznie \textbf{niezale\.zne}
od sektora grawitacyjnego: ewentualne obalenie hipotezy
o~fazowym pochodzi cechowania nie niszczy sektora grawitacyjnego.
\end{remark}

% ---------------------------------------------------------------------------
\subsection{Zasada hierarchii defekt\'ow}\label{ssec:defect-hierarchy}

\begin{axiom}[Zasada minimalno\'sci stopni swobody]%
\label{ax:minimal-dof}
Nowe stopnie swobody substratu s\k{a} dozwolone \textbf{wtedy i~tylko wtedy},
gdy s\k{a} niezb\k{e}dne do opisu klasy lokalnych defekt\'ow topologicznych,
kt\'orych nie da si\k{e} zakodowa\'c w~dotychczasowych zmiennych.
\end{axiom}

\begin{uwaga}[Hierarchia wymusze\'n]
\label{rem:defect-hierarchy-steps}
Zasada~\ref{ax:minimal-dof} generuje nast\k{e}puj\k{a}c\k{a} hierarchi\k{e}:
\begin{center}\small
\begin{tabular}{clll}
\toprule
\textbf{Krok} & \textbf{Rozszerzenie} & \textbf{Wymuszenie} & \textbf{Wynik} \\
\midrule
0 & $\hat{s}_i \in \mathbb{R}$ (modu\l{})
  & Aksjomat~A2 (substrat $\mathbb{Z}_2$)
  & pole $\Phi$, grawitacja \\[3pt]
1 & $\psi_i \in \mathbb{C}$ (faza $\theta_i$)
  & Defekty wirowe: $\oint d\theta = 2\pi n$
  & $\mathrm{U}(1)$, foton, \l{}adunek \\[3pt]
2 & $\psi_i \in \mathbb{C}^2$ (dublet)
  & Chiralno\'s\'c: $\pi_1(\mathrm{SU}(2)) = \{0\}$,
    ale $\pi_3 \neq 0$
  & $\mathrm{SU}(2) \times \mathrm{U}(1)$, bozony $W^\pm, Z$ \\[3pt]
3 & $\psi_i \in \mathbb{C}^3$ (tryplet kolorowy)
  & Konfinement: $\pi_2(\mathrm{SU}(3)/\mathrm{U}(1)^2) \neq 0$
  & $\mathrm{SU}(3)$, gluony, kwarki \\
\bottomrule
\end{tabular}
\end{center}
Ka\.zdy krok jest \textbf{wymuszony} przez istnienie klasy defekt\'ow,
kt\'orej nie mo\.zna opisa\'c na ni\.zszym poziomie.
Nie jest to ,,doklejanie'' struktur ad~hoc, lecz \textbf{dedukcyjna
sekwencja}: je\.zeli obserwujemy defekty wirowe (faz\k{e}), musimy
rozszerzy\'c substrat o~$\theta_i$; je\.zeli obserwujemy chiralno\'s\'c,
musimy przej\'s\'c do dubletu; itd.
\end{uwaga}

\begin{uwaga}[Status zasady hierarchii]
\label{rem:hierarchy-status}
Krok~0 jest \statuslabel{Twierdzenie} (sekcje~\ref{sec:pole}--\ref{sec:formalizm}).
Krok~1 jest \statuslabel{Twierdzenie} (topologiczne):
tw.~\ref{thm:D2-complex-necessary} w~dod.~\ref{app:D2-defect-hierarchy}
dowodzi, \.ze $\pi_1 \neq 0$ wymusza $M_1 = S^1$ (U(1)).
Krok~2 jest \statuslabel{Twierdzenie} (topologiczne):
tw.~\ref{thm:D2-instanton-necessary} --- $\pi_3 \neq 0$
wymusza $M_2 = S^3 \cong SU(2)$.
Krok~3 jest \statuslabel{Twierdzenie} (topologiczne + selekcyjne):
tw.~\ref{thm:D2-monopole-necessary} --- $\pi_2 \neq 0$
wymusza $SU(3)$ z~$N_c = 3$ (anulacja ABJ + AF).
Pe\l{}ny dow\'od: tw.~\ref{thm:D2-hierarchy-main}
(dod.~\ref{app:D2-defect-hierarchy}).
Zasada~\ref{ax:minimal-dof} pozostaje \statuslabel{Aksjomat}
pe\l{}no\'sci (substrat koduje wszystkie klasy defekt\'ow),
ale sama hierarchia jest teraz \statuslabel{Twierdzenie}.
Weryfikacja: \texttt{ex207\_homotopy\_defect\_chain.py} (20/20 PASS).
\end{uwaga}

% ---------------------------------------------------------------------------

\begin{proposition}[Zachowanie emergencji $\Phi$ w~pełnym substracie]\label{prop:Phi-preserved}
W~pełnym zespolonym substracie parametr porządku amplitudy:
\begin{equation}\label{eq:Phi-from-amplitude}
    \Phi(x)
    = \frac{1}{N_B}\sum_{i \in \mathcal{B}(x)}
    \langle|\psi_i|^2\rangle
    = \langle\phi_i^2\rangle_{\mathcal{B}(x)},
\end{equation}
zachowuje wszystkie własności z~sekcji~\ref{sec:pole}--\ref{sec:formalizm}.
Pole $\Phi$ nie zależy od fazy $\theta_i$.
\end{proposition}

\begin{proof}
Hamiltoniano~\eqref{eq:H-complex} jest separowalny względem amplitudy
i~fazy dla słabego sprzężenia ($J \ll m_0^2$). Część amplitudowa
identyczna z~oryginalnym substratem daje ten sam wynik coarse-grainingu.
\end{proof}

% ----------------------------------------------------------
\subsection{U(1): emergentny foton z~fazy substratu}%
\label{ssec:u1-photon}

\subsubsection{Parametr porządku fazy}

Oprócz $\Phi = \langle\phi^2\rangle$ definiujemy \textbf{blokowy
parametr fazy}:
\begin{equation}\label{eq:phase-order}
    \mathcal{A}_\mu(x) = \frac{a_{\mathrm{sub}}}{N_B}
    \sum_{(ij) \in E_\mu(\mathcal{B})}
    \mathrm{Im}\!\left[\ln\!\left(\frac{\psi_j}{\psi_i}\right)\right],
\end{equation}
gdzie suma biegnie po krawędziach $E_\mu$ prostopadłych do kierunku
$\mu$ wewnątrz bloku, a~$a_{\mathrm{sub}}$ jest krokiem kratowym.
Wielkość $\mathcal{A}_\mu$ mierzy \textbf{gradient fazowy} ---
skumulowaną różnicę faz wzdłuż kierunku $\mu$.

\begin{theorem}[Emergencja pola elektromagnetycznego]\label{thm:photon-emergence}
W~granicy ciągłej ($a_{\mathrm{sub}} \to 0$ przy ustalonej
skali fizykalnej), dla konfiguracji o~powolnie zmiennej fazie:
\begin{enumerate}[nosep]
    \item $\mathcal{A}_\mu(x) \to A_\mu(x)$ --- czteropotencjał
          elektromagnetyczny,
    \item Natężenie: $F_{\mu\nu} = \partial_\mu A_\nu
          - \partial_\nu A_\mu$ wynika z~geometrii krawędzi substratu,
    \item Działanie Maxwell'a:
          \begin{equation}\label{eq:maxwell-action}
              S_{\mathrm{EM}}
              = -\frac{1}{4\mu_0}\int F_{\mu\nu}\,F^{\mu\nu}
                \sqrt{-g}\;d^4x
          \end{equation}
          wyłania się z~części kinetycznej hamiltonianu~\eqref{eq:H-complex}
          pod podwójnym coarse-grainingiem (amplitudy i~fazy).
\end{enumerate}
\end{theorem}

\begin{proof}[Dowód (rozszerzony)]
\textbf{Krok~1: Człon kinetyczny dla sektora fazowego.}
Dla węzła $i$ substratu: $\psi_i = \phi_i\,e^{i\theta_i}$.
Człon sprzężenia sąsiada w~kierunku $\mu$:
\begin{equation}\label{eq:kinetic-step1}
    -J\,\mathrm{Re}(\psi_i^*\psi_j)
    = -J\phi_i\phi_j\cos(\theta_j - \theta_i).
\end{equation}
W~fazie uporządkowanej ($T < T_c$): $\phi_i \approx \phi_j \approx v$,
$v = \langle\phi\rangle > 0$. Dla powolnych zmian fazy:
$\theta_j - \theta_i = a_{\mathrm{sub}}\,(\partial_\mu\theta) + \mathcal{O}(a_{\mathrm{sub}}^2)$.
Rozwijamy do drugiego rzędu w~$a_{\mathrm{sub}}$:
\begin{equation}\label{eq:kinetic-step2}
    -Jv^2\cos(a_{\mathrm{sub}}\,\partial_\mu\theta)
    = -Jv^2\!\left[1 - \tfrac{1}{2}\,a_{\mathrm{sub}}^2\,(\partial_\mu\theta)^2
      + \mathcal{O}(a_{\mathrm{sub}}^4)\right].
\end{equation}
Stały człon $-Jv^2$ wchłania się do energii próżni.
Przechodząc do granicy ciągłej ($a_{\mathrm{sub}} \to 0$, przy stałej
skali fizycznej): gęstość Lagrangiana sektora fazowego wynosi
\begin{equation}\label{eq:Lphase}
    \mathcal{L}_{\mathrm{phase}}
    = \frac{Jv^2}{2}\,a_{\mathrm{sub}}^{2}\sum_{\mu = 1}^{4}(\partial_\mu\theta)^2.
\end{equation}

\textbf{Krok~2: Identyfikacja pola elektromagnetycznego.}
Definiujemy czteropotencjał przez gradient fazy:
\begin{equation}\label{eq:Amu-from-theta}
    A_\mu \;\equiv\; \frac{\hbar_0}{e}\,\partial_\mu\theta,
\end{equation}
gdzie $e$ jest ładunkiem elementarnym (stałą integracji --- patrz uwaga
o~$\alpha_{\mathrm{em}}$ poniżej).
Tensor pola $F_{\mu\nu} = \partial_\mu A_\nu - \partial_\nu A_\mu
= \frac{\hbar_0}{e}(\partial_\mu\partial_\nu\theta
- \partial_\nu\partial_\mu\theta)$.
Dla topologicznie trywialnych konfiguracji (gradient dokładny):
$F_{\mu\nu} = 0$. Dla konfiguracji z~nietrywialną topologią fazy
(wiry fazowe: $\oint\partial_\mu\theta\,dl^\mu = 2\pi n$, $n\in\mathbb{Z}$):
$F_{\mu\nu} \neq 0$. Najmniejszy obwód plakietki ($\mu\nu$-kwadrat)
daje:
\begin{equation}\label{eq:plaquette}
    \theta_\square
    = \oint_{(\mu\nu)} \partial_\lambda\theta\,dl^\lambda
    = a_{\mathrm{sub}}^2\,\frac{e}{\hbar_0}\,F_{\mu\nu}(x_\square).
\end{equation}

\textbf{Krok~3: Działanie Maxwella z~plakietek.}
Sumując człon $\sim(\partial_\mu\theta)^2$ z~równania~\eqref{eq:Lphase}
po wszystkich parach kierunków $(\mu < \nu)$ i~podstawiając
$\partial_\mu\theta = (e/\hbar_0)A_\mu$:
\begin{align}
    S_{\mathrm{phase}}
    &= \frac{Jv^2 a_{\mathrm{sub}}^2}{2}
       \int\!\sum_\mu(\partial_\mu\theta)^2\,d^4x
    \;=\; \frac{Jv^2 a_{\mathrm{sub}}^2}{2}\cdot
          \frac{e^2}{\hbar_0^2}
          \int\!\sum_\mu A_\mu A^\mu\,d^4x \label{eq:Sphase-step1}\\
    &\supset\;
    \frac{Jv^2 a_{\mathrm{sub}}^2\,e^2}{2\hbar_0^2}
    \int\!\sum_{\mu<\nu}F_{\mu\nu}^2\,d^4x, \label{eq:Sphase-maxwell}
\end{align}
gdzie $\supset$ oznacza część transwersalną (fizyczne polaryzacje).
Porównując z~działaniem Maxwella $S_{\mathrm{EM}}
= -\frac{1}{4\mu_0}\int F_{\mu\nu}F^{\mu\nu}\,d^4x$, uzyskujemy:
\begin{equation}\label{eq:mu0-substrate}
    \boxed{
      \frac{1}{\mu_0}
      = \frac{2Jv^2 a_{\mathrm{sub}}^2\,e^2}{\hbar_0^2}.
    }
\end{equation}

\textbf{Krok~4: Niezmienniczość cechowania.}
Transformacja $\theta_i \to \theta_i + \lambda(x_i)$
odpowiada $A_\mu \to A_\mu + \frac{\hbar_0}{e}\partial_\mu\lambda$ ---
standardowej transformacji cechowania $\mathrm{U}(1)$.
Tensor $F_{\mu\nu} = \partial_\mu A_\nu - \partial_\nu A_\mu$ jest
na nią \textbf{dokładnie niezmienniczy}
(weryfikacja numeryczna: test G2 w~\texttt{gauge\_emergence.py}).

\textbf{Krok~5: Bezmasowość.}
Człon~\eqref{eq:Lphase} zależy od $\theta$ wyłącznie przez
$\partial_\mu\theta$ --- brak członu $\propto\theta^2$.
Symetria $\theta_i \to \theta_i + \mathrm{const}$ jest ciągła i~nieuszkodzona.
W~przestrzeni pędowej: propagator $\propto k^{-2}$ (bezmasowy).
Foton pozostaje bezmasowy przy dowolnym obcięciu UV $\Lambda \leq \ell_P^{-1}$,
bo ochrona jest topologiczna (symetria $\mathrm{U}(1)$ substratu),
nie fine-tuningowa.
\end{proof}

\begin{remark}[Kluczowa różnica od standardowego elektromagnetyzmu]
W~standardowej elektrodynamice $A_\mu$ postulujemy jako pole na
gotowej rozmaitości. W~TGP $A_\mu$ jest \textbf{gradientem fazy
substratu} --- konstytuuje się razem z~$g_{\mu\nu}$ (przez $\Phi$),
nie na $g_{\mu\nu}$. Metryka i~pole elektromagnetyczne mają
wspólne źródło substratowe.
\end{remark}

\subsubsection{Stała struktury subtelnej}

\begin{proposition}[Stała sprzężenia elektromagnetycznego]%
\label{prop:alpha-em}
Stała struktury subtelnej $\alpha_{\mathrm{em}} = e^2/(4\pi\varepsilon_0\hbar c)$
wyraża się przez parametry substratu:
\begin{equation}\label{eq:alpha-em-substrate}
    \alpha_{\mathrm{em}}
    = \frac{J^2\,a_{\mathrm{sub}}^4\,v^4}
           {4\pi\,\hbar\,c_0\,\varepsilon_0^{-1}},
\end{equation}
gdzie $v = \langle\phi\rangle$ jest parametrem porządku amplitudy,
a~$J$ --- sprzężeniem między sąsiadami. W~granicy renormalizacji
do skali $\lP$ wartość $\alpha_{\mathrm{em}} \approx 1/137$
wyznacza relację między sprzężeniem substratowym $J$ a~skalą $v$.
\end{proposition}

\begin{remark}[Zamknięcie O12 -- przepływ RG stałej $\alpha_{\mathrm{em}}$]%
\label{rem:O12-closed}
Problem O12 jest \textbf{częściowo zamknięty} dzięki 1-pętlowemu przepływowi
RG grupy renormalizacji dla $U(1)_Y$ ze standardowymi progami SM.
Odwrotny przepływ od skali eksperymentalnej $\mu = m_e$ do skali Plancka daje:
\begin{equation}\label{eq:alpha-Planck}
    \frac{1}{\alpha_{\mathrm{em}}(\lP)}
    = \frac{1}{\alpha_{\mathrm{em}}(m_e)}
      - \frac{b_0}{3\pi}\,\ln\!\frac{\lP}{m_e}
    \approx 94.09,
    \qquad
    \alpha_{\mathrm{sub}} = \alpha_{\mathrm{em}}(\lP) \approx \frac{1}{94.09},
\end{equation}
gdzie $b_0 = \tfrac{4}{3}\sum_f Q_f^2\,N_{c,f}$ przebiega po fermionach SM.
Sprzężenie substratowe:
\begin{equation}\label{eq:J-substrate-alpha}
    J_{\mathrm{sub}}
    = \sqrt{4\pi\,\alpha_{\mathrm{sub}}}
    \approx 0.366
    \quad (\text{rząd jedności w jednostkach Plancka}).
\end{equation}
Weryfikacja numeryczna: \texttt{scripts/alpha\_em\_rg\_flow.py} -- 9/9 PASS\@.
Otwarte: pełne wielopętlowe wyprowadzenie $\alpha_{\mathrm{sub}}$ z~parametrów
hamiltonowskich substratu (eq.~\eqref{eq:alpha-em-substrate}).
\end{remark}

% ----------------------------------------------------------
\subsection{SU(2)×U(1): substrat dwuskładnikowy i~sektor
  elektrosłaby}\label{ssec:su2u1}\label{ssec:su2-su3}

\subsubsection{Rozszerzenie do dubletu}

\begin{axiom}[Substrat SU(2)-symetryczny]\label{ax:su2-substrate}
Rozszerzamy stan węzła do \textbf{dubletu}:
\begin{equation}\label{eq:doublet-state}
    \Psi_i = \begin{pmatrix} \psi_i^{(1)} \\ \psi_i^{(2)} \end{pmatrix}
    \in \mathbb{C}^2,
\end{equation}
z~hamiltonianem niezmienniczym względem globalnej $SU(2) \times U(1)$:
\begin{equation}\label{eq:H-su2}
    H_\Gamma^{(SU2)} = \sum_i \left[
        \frac{m_0^2}{2}\,|\Psi_i|^2
        + \frac{\lambda_0}{4}\,|\Psi_i|^4
    \right]
    - J\sum_{\langle ij\rangle} \mathrm{Re}
      \bigl(\Psi_i^\dagger\,\Psi_j\bigr).
\end{equation}
\end{axiom}

Parametr porządku dubletu:
\begin{equation}\label{eq:doublet-order}
    \langle\Psi_i\rangle
    = v_W\begin{pmatrix} 0 \\ 1 \end{pmatrix}
    \quad (T < T_{c,W}),
\end{equation}
gdzie $v_W$ jest odpowiednikiem vev Higgsa, a~$T_{c,W}$ temperaturą
substratową elektrosłabego łamania symetrii.

\begin{theorem}[Emergencja bosonu W, Z i~fotonu]\label{thm:su2-emergence}
Z~hamiltonianu~\eqref{eq:H-su2} w~granicy ciągłej wyłaniają się:
\begin{enumerate}[nosep]
    \item \textbf{Trzy pola $SU(2)$}: $W_\mu^a$ ($a = 1,2,3$)
          z~fluktuacji izospinorowej fazy $\Psi_i$,
    \item \textbf{Jedno pole $U(1)$}: $B_\mu$ z~globalnej
          fazy dubletu (analogicznie jak foton w~\S\ref{ssec:u1-photon}),
    \item \textbf{Trzy masywne bozony}: $W^\pm$, $Z^0$
          z~mechanizmu substratowego łamania $SU(2)_L \times U(1)_Y
          \to U(1)_{\mathrm{EM}}$ (gdy $v_W \neq 0$),
    \item \textbf{Bezmasowy foton}: z~ocalałej kombinacji
          $A_\mu = B_\mu\cos\theta_W + W_\mu^3\sin\theta_W$.
\end{enumerate}
\end{theorem}

\begin{proof}[Dowód (rozszerzony)]
\textbf{Krok~1: Parametryzacja wokół minimum.}
W~fazie $v_W \neq 0$ rozwijamy:
\begin{equation}\label{eq:su2-fluct}
    \Psi_i = \left(\frac{v_W + H_i/\sqrt{2}}{\sqrt{2}}\right)
    \exp\!\left(\frac{i\xi_i^a\sigma^a}{v_W}\right)
    \begin{pmatrix}0\\1\end{pmatrix},
\end{equation}
gdzie $H_i$ --- fluktuacja amplitudy (Higgs), $\xi_i^a$ ---
trzy fazy izospinorowe ($a=1,2,3$, $\sigma^a$ --- macierze Pauliego).

\textbf{Krok~2: Człony kinetyczne z~hamiltonianu.}
Część $-J\,\mathrm{Re}(\Psi_i^\dagger\Psi_j)$ rozwijamy tak jak
w~dowodzie twierdzenia~\ref{thm:photon-emergence} (Krok~1--2), ale
teraz dla \emph{dubletu}. Człony izospinorowe dają:
\begin{equation}\label{eq:su2-kinetic}
    \mathcal{L}_{W,\mathrm{kin}}
    = \frac{Jv_W^2 a_{\mathrm{sub}}^2}{2}
    \sum_{\mu,a}(\partial_\mu\xi^a)^2.
\end{equation}
Identyfikujemy: $W_\mu^a = \frac{\hbar_0}{g_W}\partial_\mu\xi^a$
z~$g_W^2 = 2Je^2 a_{\mathrm{sub}}^2/\hbar_0^2$ (analogicznie do~\eqref{eq:mu0-substrate}).

\textbf{Krok~3: Masy bozonów z~łamania symetrii.}
Po wybraniu vev~\eqref{eq:doublet-order} Goldstone'owie
$\xi^a$ zostają ``pochłonięte'' przez pola cechowania.
Masa wynika z~członu kinetycznego dubletu przy $W_\mu^a \neq 0$:
\begin{equation}\label{eq:W-mass-deriv}
    \mathcal{L}_{\mathrm{mass}}
    = \frac{v_W^2}{4}
      \Bigl[g_W^2(W_\mu^1{}^2 + W_\mu^2{}^2)
      + (g_W W_\mu^3 - g_B B_\mu)^2\Bigr].
\end{equation}
Stąd:
\begin{align}
    m_W^2 &= \frac{g_W^2 v_W^2}{4}, \label{eq:mW-derived}\\
    m_Z^2 &= \frac{(g_W^2 + g_B^2)\,v_W^2}{4}, \label{eq:mZ-derived}\\
    m_\gamma &= 0
    \quad\text{(kombinacja ortogonalna
    $A_\mu = B_\mu\cos\theta_W + W_\mu^3\sin\theta_W$ bezmasowa)}. \label{eq:mgamma-zero}
\end{align}
Kąt Weinsberga: $\tan\theta_W = g_B/g_W$. Relacja
$m_W/m_Z = \cos\theta_W$ wynika bezpośrednio ze struktury~\eqref{eq:W-mass-deriv}
i~jest niezależna od wartości $v_W$ --- predykcja TGP
spójna z~obserwacją ($m_W/m_Z = 0.881$, $\cos\theta_W = 0.877$,
różnica $\sim 0.5\%$ wynika z~radiacyjnych korekcji QCD).
\end{proof}

\begin{remark}[Higgs jako amplituda dubletu substratowego]%
\label{rem:higgs-substrate}
Pole Higgsa $H$ nie jest postulowane jako osobne pole ---
jest \textbf{fluktuacją amplitudy} dubletu substratowego
wokół $v_W$. Jego masa $m_H^2 = 2\lambda_0 v_W^2$
jest wyznaczona przez parametr samosprężenia $\lambda_0$
tego samego hamiltonianu. W~TGP mechanizm Higgsa
jest \emph{substratowym przejściem fazowym}, nie $ad hoc$
wyborem potencjału.
\end{remark}

\subsubsection{Łamanie SU(2)×U(1) jako granica fazowa}

\begin{proposition}[Temperatura elektrosłaba a~$\Phi$]\label{prop:EW-phase}
Substratowe przejście elektrosłabe zachodzi w~temperaturze $T_{c,W}$,
która koreluje z~wartością $\Phi$:
\begin{equation}\label{eq:EW-Phi-corr}
    T_{c,W} \;\propto\; \sqrt{\Phi / \PhiZero} \cdot T_{\mathrm{Planck}}.
\end{equation}
We wczesnym Wszechświecie ($\Phi = \PhiZero\psi$ z~$\psi_{\mathrm{ini}} = 1$)
temperatura substratowa przekraczała $T_{c,W}$, więc $SU(2) \times U(1)$
było niezłamane. Łamanie następuje naturalnie w~epoce
$T_{\mathrm{cosm}} \sim v_W \approx 246\;\mathrm{GeV}$,
spójnie z~kosmologią standardową.
\end{proposition}

\subsubsection{$SU(2)$ jako podw\'ojne pokrycie $SO(3)$
  chiralnego substratu}\label{sssec:su2-chiral}

\begin{proposition}[$SU(2)_L$ z~chiralnej asymetrii substratu]%
\label{prop:su2-from-chirality}
\statuslabel{Propozycja}
Emergencja $SU(2)_L$ (a nie $SU(2)_R$) z~substratu ma
\textbf{geometryczne} \'zr\'od\l{}o w~chiralnej asymetrii
potencja\l{}u $V(g) = g^3/3 - g^4/4$ (dod.~\ref{app:K2-chiralnosc}).

\medskip\noindent\textbf{Argument w~trzech krokach.}

\emph{Krok~1: Dwa sektory chiralne substratu.}\;
Potencja\l{} $V(g)$ ma $V'''(1) = -4 \neq 0$ (prop.~\ref{prop:K2-Vasym}),
co \l{}amie symetri\k{e} $g \leftrightarrow 2-g$ wok\'o\l{} pr\'o\.zni.
Pary chiralne soliton\'ow $(g_0^R, g_0^L)$ maj\k{a} r\'o\.zne amplitudy
ogonowe: $A_{\rm tail}^R \neq A_{\rm tail}^L$, wi\k{e}c r\'o\.zne masy
(tw.~\ref{thm:K2-chiral-selection}). Stosunek $m_R/m_L \approx 2{,}03$
(dod.~\ref{app:K2-chiralnosc}).

\emph{Krok~2: Przestrze\'n orbit soliton\'ow.}\;
Konfiguracja solitonowa jest sferyczna: profil $g(r)$ okre\'sla punkt
w~$SO(3)$ (orientacja rdzenia). Dwie ga\l{}\k{e}zie chiralne
$(g_0^R, g_0^L)$ definiuj\k{a} \textbf{podw\'ojne pokrycie}:
\begin{equation}\label{eq:su2-cover}
    \widetilde{\mathcal{C}}_{\rm sol} \simeq SU(2) \;\cong\; S^3
    \;\xrightarrow{\;2:1\;}\; SO(3) \;\simeq\; \mathcal{C}_{\rm sol},
\end{equation}
gdzie ka\.zdy punkt $SO(3)$ (orientacja) ma dwa podniesienia:
lewe ($g_0^L < g^*$) i~prawe ($g_0^R > g^*$), rozdzielone
barier\k{a} duchow\k{a} $g^* = e^{-1/4}$.%
\footnote{W~sformu\l{}owaniu kanonicznym (Formulacja~A): $K(g) = g^4$,
punkt duchowy nie wyst\k{e}puje ($K(g) > 0$ dla $g > 0$).
Bariera chiralna wynika z~topologii $V(g)$, nie z~$K(g) = 0$.
$g_0^e = 0{,}86941$, czarna dziura: $g \to 0$.
Patrz: \texttt{tgp\_companion.tex}.}
To jest dok\l{}adnie topologia $SU(2)$, nie postulat.

\emph{Krok~3: Selekcja chiralna.}\;
Bariera duchowa $K_{\rm sub}(g^*) > 0$ (tw.~\ref{thm:ghost-free-soliton})
blokuje g\l{}\k{e}bokie kinki --- ga\l{}\k{a}\'z $L$ jest
energetycznie dost\k{e}pna tylko dla lekkich soliton\'ow
($g_0^L \lesssim g^*$). W~efekcie sektor elektros\l{}aby
dzia\l{}a \textbf{wy\l{}\k{a}cznie na lewor\k{e}cznych}:
$SU(2)_L$ emerguje naturalnie, $SU(2)_R$ jest energetycznie
zablokowana.
\end{proposition}

\begin{remark}[Warunek falsyfikowalno\'sci sektora $SU(2)$]%
\label{rem:su2-falsify}
Propozycja~\ref{prop:su2-from-chirality} daje
\textbf{dwie mierzalne predykcje}:
\begin{enumerate}[nosep]
  \item Stosunek $m_W/m_Z = \cos\theta_W$ wynika ze~struktury
    dubletu~\eqref{eq:W-mass-deriv} \textbf{bez wolnych parametr\'ow}
    (obserwacja: $0{,}881$ vs.~$\cos\theta_W = 0{,}877$; zgodno\'s\'c $0{,}5\%$,
    r\'o\.znica tłumaczona korekcjami radiacyjnymi QCD).
  \item Chiralny stosunek mas $m_R/m_L \approx 2{,}03$ (z~asymetrii $V(g)$)
    --- je\'sli przysz\l{}e pomiary neutrin prawor\k{e}cznych
    poka\.z\k{a} $m_{\nu_R}/m_{\nu_L} \neq 2 \pm 20\%$,
    mechanizm chiralny TGP jest obalony.
\end{enumerate}
\end{remark}

% ----------------------------------------------------------
\subsection{SU(3): substrat trójskładnikowy i~sektor koloru}%
\label{ssec:su3-color}

\subsubsection{Substrat o~lokalnej symetrii $\mathbb{Z}_3$}

Substrat dwój\-kowy (odpowiedzialny za elektrosłabe) rozszerzamy
do \textbf{trójskładnikowego}:

\begin{axiom}[Substrat kolorowy]\label{ax:color-substrate}
Każdy węzeł $i \in V$ nosi \textbf{trójskładnikowy} stan:
\begin{equation}\label{eq:triplet-state}
    \Xi_i = \begin{pmatrix} \psi_i^r \\ \psi_i^g \\ \psi_i^b \end{pmatrix}
    \in \mathbb{C}^3,
\end{equation}
odpowiadający trzem \emph{kolorom substratowym}
($r$, $g$, $b$). Hamiltoniano jest niezmienniczy względem globalnej
$SU(3)_c$:
\begin{equation}\label{eq:H-color}
    H_\Gamma^{(SU3)} = \sum_i \left[
        \frac{m_0^2}{2}\,|\Xi_i|^2
        + \frac{\lambda_0}{4}\,|\Xi_i|^4
    \right]
    - J_c\sum_{\langle ij\rangle} \mathrm{Re}
      \bigl(\Xi_i^\dagger\,\Xi_j\bigr).
\end{equation}
\end{axiom}

Kluczowe: substrat \textbf{nie} jest siatką barwną w~zwykłym
sensie --- kolory są wewnętrznymi stopniami swobody węzłów substratu,
nie etykietami w~przestrzeni.

\begin{theorem}[Emergencja gluonów]\label{thm:gluon-emergence}
Z~hamiltonianu~\eqref{eq:H-color} w~granicy ciągłej wyłaniają się:
\begin{enumerate}[nosep]
    \item \textbf{Osiem pól gluonów} $G_\mu^a$ ($a = 1,\ldots,8$)
          z~gradientów fazy tripoletu $\Xi_i$,
    \item \textbf{Działanie Yang--Millsa}:
          $S_{SU(3)} = -\frac{1}{2g_s^2}\int
          \mathrm{Tr}[F_{\mu\nu}F^{\mu\nu}]\sqrt{-g}\,d^4x$,
    \item \textbf{Potencjał konfinowania} gluonów: wynika
          bezpośrednio z~re\.zimu III TGP (studnia przestrzenna,
          \S\ref{ssec:studnia}), który działa na kwarki jako
          nośniki ładunku kolorowego, nie tylko grawitacyjnego.
\end{enumerate}
\end{theorem}

\begin{proof}[Dowód (rozszerzony)]
\textbf{Krok~1: Stopnie swobody fazy tripletu.}
Parametryzacja: $\Xi_i = |\Xi_i|\,U_i$, gdzie $U_i \in SU(3)/U(1)^2$.
Globalną fazę $U(1)_B$ (barion) pomijamy. Stopnie swobody
fazowe $SU(3)$ to \textbf{osiem} niezależnych parametrów
(wymiar algebry $\mathfrak{su}(3) = 8$).

\textbf{Krok~2: Osiem pól gluonowych.}
Dla krawędzi w~kierunku $\mu$ rozwijamy sprzężenie do drugiego rzędu
w~gradientach (analogicznie do kroku~1 twierdzenia~\ref{thm:photon-emergence}):
\[
  -J_c\,\mathrm{Re}(\Xi_i^\dagger\Xi_j)
  = -J_c v_c^2 + \frac{J_c v_c^2 a_{\mathrm{sub}}^2}{2}
    \sum_{a=1}^{8}(\partial_\mu\xi^a)^2 + \ldots
\]
gdzie $\xi^a$ są parametrami Eulera rozkładu $U_j = U_i\exp(ia_{\mathrm{sub}}\xi^a\lambda^a/2)$,
$\lambda^a$ --- macierze Gell-Manna. Identyfikujemy:
\begin{equation}\label{eq:gluon-id}
    G_\mu^a = \frac{\hbar_0}{g_s}\,\partial_\mu\xi^a,
    \qquad
    g_s^2 = \frac{2J_c e^2 a_{\mathrm{sub}}^2}{\hbar_0^2}.
\end{equation}

\textbf{Krok~3: Nieabelowy tensor pola.}
W~odróżnieniu od U(1), pola $G_\mu^a$ mają samosprzężenie
ponieważ $[\lambda^a, \lambda^b] \neq 0$ ($f^{abc} \neq 0$).
Dla nie-abelowego sprzężenia krawędź--krawędź~w~plakiecie daje:
\begin{equation}\label{eq:SU3-plaq}
    F_{\mu\nu}^a = \partial_\mu G_\nu^a - \partial_\nu G_\mu^a
    + g_s f^{abc} G_\mu^b G_\nu^c,
\end{equation}
gdzie $f^{abc}$ --- stałe struktury $SU(3)$.
Człon $g_s f^{abc}G_\mu^b G_\nu^c$ wynika z~kwadratowych poprawek
do plakiety substratowej ($\sim a_{\mathrm{sub}}^4$ w~granicy
ciągłej), koniecznych do zachowania niezmienniczości cechowania.

\textbf{Krok~4: Działanie Yang--Millsa.}
Sumując po plakietach i~przechodząc do granicy ciągłej:
\begin{equation}\label{eq:YM-from-substrate}
    S_{SU(3)}
    = -\frac{1}{2g_s^2}\int
      \mathrm{Tr}\bigl[F_{\mu\nu}F^{\mu\nu}\bigr]\sqrt{-g}\;d^4x,
\end{equation}
co jest standardowym działaniem Yanga--Millsa dla $SU(3)$.
Weryfikacja numeryczna: skrypt \texttt{gauge\_emergence.py},
sekcja~D (testy D1--D4).
\end{proof}

\subsubsection{Konfinowanie kwarków w~TGP}

\begin{remark}[Mechanizm konfinowania bez dodatkowych założeń]%
\label{rem:confinement}
W~standardowym QCD konfinowanie jest efektem nieperturbacyjnym,
trudnym do wyprowadzenia z~pierwszych zasad. W~TGP mechanizm
jest \textbf{wbudowany ontologicznie}:

\begin{itemize}[nosep]
    \item Kwarki są nośnikami ładunku kolorowego $SU(3)_c$
          --- oznacza to, że stany węzłów substratu wokół kwarku
          mają niezerowy ładunek kolorowy.
    \item Studnia przestrzenna (reżim III, \S\ref{ssec:studnia})
          działa na każdy ładunek substratowy, nie tylko
          grawitacyjny. Dla ładunków kolorowych studnia
          ma głębokość $\sim \Lambda_{\mathrm{QCD}}$.
    \item \textbf{String napięcia}: konfiguracja pola $\Xi_i$
          między kwarkiem a~antykwarkiem tworzy \emph{rurę
          kolorową} --- nitkę, wzdłuż której $|\Xi|$ jest
          podwyższone ponad wartość próżniową. Energia tej
          nitki jest proporcjonalna do długości: $V(r) \approx \sigma r$,
          co jest konfinowaniem liniowym.
\end{itemize}

Napięcie stringa $\sigma$ jest wyznaczone przez energię
konfiguracji substratowej:
\begin{equation}\label{eq:string-tension}
    \sigma = \frac{\Phi_{\mathrm{string}}}{\PhiZero} \cdot
    \frac{\gamma\,\PhiZero}{\lP^2} \cdot a_{\mathrm{sub}},
\end{equation}
gdzie $\Phi_{\mathrm{string}}$ jest wartością pola przestrzenności
wewnątrz rurki kolorowej. Numerycznie $\sigma \approx 0.18\;\mathrm{GeV}^2$
wyznacza relację między $\gamma$, $\PhiZero$ i~$a_{\mathrm{sub}}$.
\end{remark}

\begin{proposition}[Ilościowe oszacowanie napięcia stringa z~parametrów TGP]%
\label{prop:sigma-estimate}
Przy założeniach:
\begin{enumerate}[nosep]
    \item[(S1)] Rurka kolorowa ma przekrój poprzeczny $\sim a_{\mathrm{sub}}^2$
          (minimalna skala substratowa),
    \item[(S2)] Wewnątrz rurki pole $\Phi$ jest stłumione do wartości
          $\Phi_{\mathrm{string}} \approx f_{\mathrm{col}} \cdot \PhiZero$
          z~$f_{\mathrm{col}} \equiv \Phi_{\mathrm{string}}/\PhiZero
          \in (0,1)$ (reżim~III --- studnia),
    \item[(S3)] Gęstość energii rurki jest zdominowana przez gradient
          $\Phi$ na brzegu: $\varepsilon_{\mathrm{tube}}
          \sim K(\Phi_0)\,(\Delta\Phi)^2 / a_{\mathrm{sub}}^2$,
\end{enumerate}
napięcie stringa wynosi:
\begin{equation}\label{eq:sigma-from-TGP}
    \sigma \;=\; \varepsilon_{\mathrm{tube}} \cdot a_{\mathrm{sub}}^2
    \;=\; K_{\mathrm{geo}}\,\PhiZero^4 \cdot
    \frac{(1 - f_{\mathrm{col}})^2 \cdot \PhiZero^2}{a_{\mathrm{sub}}^2}
    \cdot a_{\mathrm{sub}}^2
    \;=\; K_{\mathrm{geo}}\,(1 - f_{\mathrm{col}})^2\,\PhiZero^6,
\end{equation}
gdzie użyto $K(\phi) = K_{\mathrm{geo}}\,\phi^4$
(twierdzenie~\ref{thm:D-uniqueness}),
$\Delta\Phi = \PhiZero(1 - f_{\mathrm{col}})$,
a~$K_{\mathrm{geo}} = 3/(4\PhiZero) = \kappa$ (sek.~\ref{ssec:unified-action}).

\noindent\textbf{Oszacowanie numeryczne.}
W~jednostkach Plancka ($\lP = 1$, $\PhiZero = 24{,}66$,
$a_{\Gamma} = 0{,}040$):
\begin{itemize}[nosep]
    \item $K_{\mathrm{geo}} = \kappa = 3/(4\PhiZero) = 0{,}0304$,
    \item $\sigma = \kappa\,(1-f_{\mathrm{col}})^2\,\PhiZero^6$
          (w~jednostkach Plancka, $M_{\mathrm{Pl}}^2 = 1$).
\end{itemize}
Obserwacyjna wartość $\sigma_{\mathrm{obs}} \approx 0{,}18\;\mathrm{GeV}^2
= 0{,}18/(1{,}22\times 10^{19})^2\;\mathrm{M_{Pl}^2}
\approx 1{,}2 \times 10^{-38}\;\mathrm{M_{Pl}^2}$.

Z~drugiej strony:
$\kappa\,\PhiZero^6 = 0{,}0304 \times (24{,}66)^6 \approx 6{,}8 \times 10^{6}$,
więc warunek zgodności wymaga:
\begin{equation}\label{eq:fcol-from-sigma}
    (1 - f_{\mathrm{col}})^2 = \frac{\sigma_{\mathrm{obs}}}{\kappa\,\PhiZero^6}
    \approx \frac{1{,}2 \times 10^{-38}}{6{,}8 \times 10^{6}}
    = 1{,}8 \times 10^{-45},
    \quad\Rightarrow\quad
    1 - f_{\mathrm{col}} \approx 4{,}2 \times 10^{-23}.
\end{equation}
Interpretacja: $\Phi_{\mathrm{string}} \approx \PhiZero\,(1 - 4{,}2\times 10^{-23})$
--- rurka kolorowa prawie nie zmienia pola $\Phi$.
To jest spójne z~obserwacją, że QCD działa w~reżimie
$\Phi \approx \PhiZero$ (odchylenie od próżni jest $\ll 1$).
\end{proposition}

\begin{remark}[Alternatywne oszacowanie przez $\Lambda_{\mathrm{QCD}}$]%
\label{rem:sigma-alt}
Niezależne oszacowanie: $\sigma \approx \Lambda_{\mathrm{QCD}}^2$
z~wymiaru, gdzie $\Lambda_{\mathrm{QCD}} \approx 0{,}3\;\mathrm{GeV}$
daje $\sigma \sim 0{,}09\;\mathrm{GeV}^2$ (rzędu wielkości).
W~TGP skala $\Lambda_{\mathrm{QCD}}$ powinna wynikać z~RG biegu
$\alpha_s(\Phi)$ (prop.~\ref{prop:running-coupling}):
\begin{equation}\label{eq:LQCD-from-TGP}
    \Lambda_{\mathrm{QCD}} \;\sim\; M_{\mathrm{Pl}}\,
    \exp\!\Bigl(-\frac{2\pi}{b_0\,\alpha_s(\PhiZero)}\Bigr),
\end{equation}
gdzie $b_0 = (11N_c - 2N_f)/(12\pi) = (33-6)/(12\pi) = 9/(4\pi) \approx 0{,}716$.
Odtworzenie $\Lambda_{\mathrm{QCD}} \approx 0{,}3\;\mathrm{GeV}$
wymaga:
\begin{equation}\label{eq:alphas-required}
    \alpha_s(\PhiZero) = \frac{2\pi}{b_0 \cdot \ln(M_{\mathrm{Pl}}/\Lambda_{\mathrm{QCD}})}
    = \frac{2\pi}{0{,}716 \cdot 45{,}15} \approx 0{,}194,
\end{equation}
co jest $\sim 5\times$ większe od $a_{\Gamma} = 0{,}040$.
Zatem $\alpha_s(\PhiZero) \neq a_{\Gamma}$ przy standardowym
jednopętlowym biegu QCD. Identyfikacja $\alpha_s$ z~parametrem
substratowym wymaga albo niestandardowej funkcji~$\beta$ w~TGP,
albo porzucenia tej ścieżki (problem otwarty~O15, rozwiązany
\textbf{negatywnie} przy standardowym RG).
\end{remark}

% ----------------------------------------------------------
\subsection{Pełna struktura cechowania: SM z~substratu}%
\label{ssec:full-gauge}

Łącząc sekcje~\ref{ssec:u1-photon}--\ref{ssec:su3-color},
substrat pełny TGP ma strukturę:

\begin{definition}[Pełny substrat TGP]\label{def:full-substrate}
Węzeł substratu nosi stan:
\begin{equation}\label{eq:full-substrate-state}
    \Sigma_i = (\phi_i,\;\theta_i^{\mathrm{EM}},\;\Psi_i^{\mathrm{EW}},
                \;\Xi_i^{\mathrm{QCD}}),
\end{equation}
gdzie:
\begin{itemize}[nosep]
    \item $\phi_i$ --- amplituda (generuje $\Phi$, czyli przestrze\'n),
    \item $\theta_i^{\mathrm{EM}}$ --- faza globalna $U(1)_{\mathrm{EM}}$
          (generuje foton),
    \item $\Psi_i^{\mathrm{EW}} \in \mathbb{C}^2$ --- dublet
          elektrosłaby (generuje $W^\pm$, $Z$, Higgs),
    \item $\Xi_i^{\mathrm{QCD}} \in \mathbb{C}^3$ --- triplet
          kolorowy (generuje gluony).
\end{itemize}
Całkowita symetria substratowa: $U(1)_{\mathrm{EM}} \times SU(2)_L
\times SU(3)_c$, odpowiadająca grupie cechowania Modelu Standardowego.
\end{definition}

\begin{remark}[Co TGP dodaje a~co zachowuje]
TGP \textbf{nie} modyfikuje algebry cechowania $U(1) \times SU(2)
\times SU(3)$ --- jest ona określona przez strukturę wewnętrzną
substratu. TGP \textbf{wyjaśnia}, \emph{dlaczego} ta algebra
i~żadna inna: bo substrat ma minimalne stopnie swobody kompatybilne
z~wymiarem 3+1 (wyprowadzone w~\S\ref{ssec:wymiar})
i~koniecznością konfinowania kolorów przy zachowaniu stabilności
próżni (warunek $\beta = \gamma$ przenosi się na sektor kolorowy).
\end{remark}

\begin{theorem}[Jedyność grupy cechowania SM w~3+1 wymiarach]%
\label{thm:gauge-uniqueness}
Niech substrat $\Gamma$ spełnia:
\begin{enumerate}[nosep]
    \item[(G1)] Wymiar $d_{\mathrm{space}} = 3$
          (prop.~\ref{prop:wymiar}, \S\ref{ssec:wymiar}),
    \item[(G2)] Stabilność próżni: $\beta = \gamma$ (stw.~\ref{prop:vacuum-condition}),
    \item[(G3)] Asymptotyczna swoboda sektora kolorowego
          (wymagana przez spójność perturbatywną w~reżimie~III),
    \item[(G4)] Zachowanie parzystości ładunkowej
          (C-parzystość: warunek minimalności chiralna).
\end{enumerate}
Wtedy minimalna kompatybilna z~tymi warunkami grupa substratowych
symetrii wewnętrznych to:
\begin{equation}\label{eq:gauge-group-unique}
    \boxed{
        G_{\mathrm{substrat}} = U(1) \times SU(2) \times SU(3).
    }
\end{equation}
\end{theorem}

\begin{proof}[Dowód (rozszerzony)]
\textbf{Krok~A: Warunek asymptotycznej swobody (G3).}
W~4D kwantowej teorii pola kolorowej, funkcja beta dla $SU(N_c)$
z~$N_f$ generacjami kwarków:
\begin{equation}\label{eq:beta-func}
    \beta(g_s) = -\frac{g_s^3}{16\pi^2}
    \left(\frac{11N_c}{3} - \frac{2N_f}{3}\right) + O(g_s^5).
\end{equation}
Asymptotyczna swoboda: $\beta < 0$, co wymaga
$11N_c > 2N_f$, czyli $N_c > 2N_f/11$.
Dla $N_f = 3$ generacji (z~additionF): $N_c > 6/11 \approx 0.55$.
Warunek ten spełniają $N_c = 1, 2, 3, \ldots$. Wykluczamy $N_c = 1$
(Abelowe $U(1)_c$ nie ma asymptotycznej swobody --- stała maleje ku
UV), $N_c = 2$ --- patrz krok~C.

\textbf{Krok~B: Anulowanie anomalii chiralnych (G4).}
Warunek Adlera--Bella--Jackiwa dla teorii bez anomalii:
$\mathrm{Tr}[Y_f^3] = 0$, gdzie $Y_f$ s\k{a} hiperladunkami.
Dla jednej generacji (lewor\k{e}czne fermiony Weyla,
prawor\k{e}czne jako sprz\k{e}\.zone):
\begin{equation}\label{eq:ABJ-cancel}
  \mathrm{Tr}[Y^3] =
  \underbrace{N_c \cdot 2 \cdot \bigl(\tfrac{1}{6}\bigr)^{\!3}}_{Q_L}
  + \underbrace{N_c \cdot \bigl(-\tfrac{2}{3}\bigr)^{\!3}}_{u_R^c}
  + \underbrace{N_c \cdot \bigl(\tfrac{1}{3}\bigr)^{\!3}}_{d_R^c}
  + \underbrace{2 \cdot \bigl(-\tfrac{1}{2}\bigr)^{\!3}}_{L_L}
  + \underbrace{\bigl(1\bigr)^3}_{e_R^c}
  = \frac{3 - N_c}{4}.
\end{equation}
Warunek $\mathrm{Tr}[Y^3] = 0$ daje:
\begin{equation}\label{eq:Nc-from-ABJ}
    \boxed{N_c = 3 \quad\text{(dok\l{}adnie).}}
\end{equation}
Jest to jedyne rozwi\k{a}zanie ca\l{}kowite --- i~jest \textbf{dok\l{}adne},
nie przybli\.zone.
Dla $N_c = 2$: $\mathrm{Tr}[Y^3] = 1/4 \neq 0$.
Dla $N_c = 4$: anomalia $= -1 + 4\cdot(7/27) = +1/27 \neq 0$.

Wniosek: \textbf{$N_c = 3$ jest jedyną wartością całkowitą}, przy której
anomalie chiralne się znoszą przy ładunkach SM, co jest wymagane
przez spójność teorii (G4). Uwaga: warunek~\eqref{eq:Nc-from-ABJ}
zakłada ładunki SM jako dane --- dynamiczne wyprowadzenie ładunków
z~substratu pozostaje problemem otwartym (O13).

\textbf{Krok~C: Wykluczenie $N_c = 2$.}
$SU(2)$ nie pozwala na konfinowanie liniowe (brak konfiguracji
topologicznych z~napięciem stringa) w~3+1D --- co sprzeciwia się
reżimowi~III TGP. Ponadto anomalia chiralna dla $N_c = 2$
nie zeruje się (krok~B).

\textbf{Krok~D: Wykluczenie grup wyjątkowych.}
Grupy $G_2$, $F_4$, $E_6$, $E_7$, $E_8$ nie mają naturalnych
warunków próżniowych ($\beta = \gamma$ dla sektora $\Phi$) ---
ich centra $Z(G_2) = \{e\}$ (trywialne) uniemo\.zliwiaj\k{a}
sp\'ojne przeniesienie warunku $\mathbb{Z}_2$ substratu na sektor kolorowy.

\textbf{Krok~E: Sektor elektrosłaby.}
Minimalna grupa dająca masywne $W^\pm$, $Z$ przy bezmasowym~$\gamma$
to $SU(2) \times U(1)$ z~łamaniem $\to U(1)_{\mathrm{EM}}$.
Każda mniejsza (tylko $U(1)$) nie daje prądów naładowanych ($W^\pm$);
każda większa ($SU(3)_W$ itp.) daje dodatkowe niezaobserwowane bozony.
Minimalność (G2) selekcjonuje $SU(2) \times U(1)$.
\end{proof}

\begin{remark}[Status twierdzenia~\ref{thm:gauge-uniqueness}]%
\label{rem:gauge-uniqueness-status}
Twierdzenie~\ref{thm:gauge-uniqueness} ma charakter
\textbf{selekcji warunkami spójności}, nie ścisłego dowodu
matematycznego. Niemniej argument wskazuje,
że struktura SM jest \emph{wyjątkowa} w~ramach TGP,
nie arbitralna.
\end{remark}

\begin{theorem}[Odporno\'s\'c anulacji anomalii na dynamiczne $\Phi$ (O13)]%
\label{thm:anomaly-Phi-robust}
\statuslabel{Twierdzenie (topologiczne)}

Warunek anulacji anomalii chiralnej ABJ:
\begin{equation}\label{eq:anomaly-ABJ-Phi}
  \mathcal{A}[U(1)_Y^3]
  = \sum_f Y_f^3
  = \tfrac{3}{4} - \tfrac{N_c}{4}
  = 0
  \quad\Longleftrightarrow\quad N_c = 3
\end{equation}
jest \textbf{dok\l{}adnie niezmienniczy} wzgl\k{e}dem dynamiki
pola $\Phi$ (i~wynikaj\k{a}cej metryki $g_{\mu\nu}(\Phi)$).
\end{theorem}

\begin{proof}
Dow\'od opiera si\k{e} na trzech niezale\.znych argumentach,
z~kt\'orych ka\.zdy sam wystarcza:

\medskip\noindent
\textbf{Argument~I: Indeks Atiyaha--Singera.}
Anomalia chiralna jest \textbf{indeksem} operatora Diraca
sprz\k{e}\.zonego z~polem cechowania na rozmaitosc
z~metryk\k{a} $g_{\mu\nu}$:
\begin{equation}\label{eq:AS-index}
  \mathcal{A} = \mathrm{ind}(\slashed{D})
  = \int \mathrm{ch}(F) \wedge \hat{A}(R),
\end{equation}
gdzie $\mathrm{ch}(F)$ jest charakterem Cherna wi\k{a}zki
cechowania, a~$\hat{A}(R)$ --- genus Dirac rozmaitosc
z~krzywizn\k{a} $R$.
Kluczowa w\l{}asno\'s\'c: indeks jest \textbf{inwarianten
topologicznym} --- zale\.zy wy\l{}\k{a}cznie od klasy
Cherna wi\k{a}zki, nie od metryki.
W~TGP: $g_{\mu\nu} = g_{\mu\nu}(\Phi)$ zmienia si\k{e}
dynamicznie, ale \textbf{klasa Cherna} pola cechowania
jest zdeterminowana przez topologi\k{e} substratu
(klas\k{e} homotopii $\pi_k$), nie przez $\Phi$.
Wi\k{e}c $\mathrm{ind}(\slashed{D})$ nie zale\.zy od $\Phi$.

\medskip\noindent
\textbf{Argument~II: Topologiczny charakter \l{}adunk\'ow.}
Hiperladunki $Y_f$ s\k{a} kwantami topologicznymi substratu
(prop.~\ref{prop:charge-quantization-gauge}), zdeterminowanymi
przez klas\k{e} homotopii $\pi_1$ p\k{e}tli substratowych.
Dynamika $\Phi$ modyfikuje metryk\k{e} (d\l{}ugo\'sci, czasy),
ale \textbf{nie topologi\k{e}} --- $Y_f$ s\k{a} sta\l{}ymi
wymiernymi, niezale\.znymi od~$\Phi$.
Anomalia~\eqref{eq:anomaly-ABJ-Phi} jest \textbf{wielomianem
sto\l{}ej} w~hiperladunkach:
$\mathcal{A} = (3 - N_c)/4$ nie zawiera $\Phi$.

\medskip\noindent
\textbf{Argument~III: Twierdzenie Adlera--Bardeena.}
Anomalia jest \textbf{jednoobiegowo dok\l{}adna}:
$\mathcal{A}$ nie otrzymuje poprawek od wy\.zszych p\k{e}tli
(tw.~Adlera--Bardeena, 1969).
Pole $\Phi$ modyfikuje propagatory (przez $c(\Phi)$,
$\hbar(\Phi)$), ale \textbf{nie zmienia struktury
jednoobiegowej} --- diagramy tr\'ojk\k{a}tne maj\k{a}
t\k{e} sam\k{a} topologi\k{e} (ten sam indeks!) niezale\.znie
od $\Phi$.

\medskip\noindent
\textbf{Weryfikacja numeryczna:}
Skrypt \texttt{anomaly\_cancellation.py} (20/20~PASS)
i~\texttt{ex208\_anomaly\_ewsb\_validation.py}
potwierdzaj\k{a} anulacj\k{e} dla $N_c = 3$ i~jej brak
dla $N_c \neq 3$, niezale\.znie od profilu $\Phi(r)$.
\end{proof}

\begin{remark}[Status O13]\label{rem:O13-closed}
Problem O13 jest \textbf{zamkni\k{e}ty}: awans z~Otwarty
do~\statuslabel{Twierdzenie}.
Dow\'od korzysta z~twierdzenia indeksowego Atiyaha--Singera
(ustalony wynik matematyczny), twierdzenia Adlera--Bardeena
(ustalony wynik QFT) i~topologicznego charakteru
\l{}adunk\'ow w~TGP.
\end{remark}

\begin{proposition}[Kwantyzacja ładunku z~topologii substratu]%
\label{prop:charge-quantization-gauge}
W~TGP kwantyzacja ładunku elektrycznego wynika z~dwóch niezależnych
mechanizmów:
\begin{enumerate}[nosep]
    \item \textbf{Topologiczny}: Faza substratowa $\theta_i^{\mathrm{EM}}$
          jest zwarta ($\theta \in [0, 2\pi)$), więc ładunek
          $Q = (1/2\pi)\oint \mathrm{d}\theta$ jest całkowity
          w~jednostkach $e$ (argument Diraca na substracie).
    \item \textbf{Anomalijny}: Warunek anulowania anomalii~\eqref{eq:ABJ-cancel}
          z~$N_c = 3$ wymaga ładunków ułamkowych $Q_u = 2/3$,
          $Q_d = -1/3$ --- ale te są wielokrotnościami $e/3$,
          a~\textbf{konfinowanie} gwarantuje, że obserwowalne
          stany (hadrony) mają ładunek całkowity.
\end{enumerate}
Łącząc oba mechanizmy: ładunek jest skwantyzowany w~jednostkach $e/3$
fundamentalnie i~w~jednostkach $e$ obserwacyjnie (konfinowanie).
Jest to \textbf{predykcja TGP}: nie istnieją wolne ładunki ułamkowe
(warunek obalenia~K8 z~\S\ref{sec:predykcje}).
\end{proposition}

% ----------------------------------------------------------
\subsection{Sprzężenie pól cechowania z~$\Phi$}%
\label{ssec:gauge-Phi-coupling}

Pola cechowania, wyłaniając się z~tego samego substratu co $\Phi$,
są sprzężone z~$\Phi$ w~sposób nietrywialny:

\begin{proposition}[Efektywna stała sprzężenia zależy od $\Phi$]%
\label{prop:running-coupling}
Stałe sprzężenia elektromagnetycznego i~silnego zależą od lokalnego
$\Phi$ przez renormalizację substratową:
\begin{align}
    \alpha_{\mathrm{em}}(\Phi)
    &= \alpha_{\mathrm{em},0}
       \left(\frac{\Phi}{\PhiZero}\right)^{\delta_{\mathrm{em}}},
    \label{eq:alpha-Phi} \\
    \alpha_s(\Phi)
    &= \alpha_{s,0}
       \left(\frac{\Phi}{\PhiZero}\right)^{-\delta_s},
    \label{eq:alphas-Phi}
\end{align}
gdzie $\delta_{\mathrm{em}} > 0$ (sprzężenie EM rośnie w~gęstej przestrzeni)
i~$\delta_s > 0$ (sprzężenie silne maleje --- asymptotyczna swoboda
jest wzmocniona przez gęstość przestrzeni).
\end{proposition}

\begin{remark}[Testowalność]\label{rem:coupling-test}
Zależność~\eqref{eq:alpha-Phi}--\eqref{eq:alphas-Phi} jest
\textbf{predykcją TGP}: stałe sprzężenia powinny zależeć od
lokalnego potencjału $\Phi$ (czyli od pola grawitacyjnego).
Efekt jest rzędu $\delta_{\mathrm{em}} \cdot U_N \sim 10^{-10}$
w~Układzie Słonecznym (niemierzalny), ale może być istotny
w~silnych polach (blisko gwiazd neutronowych, CD).
\end{remark}

% ----------------------------------------------------------
\subsection{Most masa--sprz\k{e}\.zenie: $\alpha_s$ z~sektora solitonowego}%
\label{ssec:alphas-soliton-bridge}
% ----------------------------------------------------------

\begin{proposition}[$\alpha_s(M_Z)$ z~$\varphi$-FP i~$\Phi_0$]%
\label{prop:alphas-phiFP}
Sta\l{}a sprz\k{e}\.zenia silnego w~skali $M_Z$ jest wyra\.zalna
jako iloczyn trzech czynnik\'ow --- kolorowego, solitonowego
i~grawitacyjnego:
\begin{equation}\label{eq:alphas-formula}
  \boxed{\alpha_s(M_Z)
    = \underbrace{\frac{N_c^2}{2}}_{\text{kolor: } T_F \cdot N_c}
    \;\cdot\;
    \underbrace{g_0^{\,e}}_{\text{soliton: }\varphi\text{-FP}}
    \;\cdot\;
    \underbrace{\frac{1}{\PhiZero}}_{\text{grawit.: }\sim\kappa}
    \;=\; \frac{N_c^3\,g_0^{\,e}}{8\,\PhiZero},}
\end{equation}
gdzie $g_0^{\,e} = 0.8694$ to warto\'s\'c $\varphi$-FP
z~ODE substratowego (parametryzacja $\alpha = 1$,
$K_{\rm sub}(g) = g^2$), $N_c = 3$, $\PhiZero = 36\Omega_\Lambda$.

\textbf{Interpretacja czynnik\'ow:}
\begin{itemize}[nosep]
  \item $T_F \cdot N_c = N_c^2/2$: g\k{e}sto\'s\'c sprz\k{e}\.zenia
    soliton--substrat w~przestrzeni kolorowej
    ($T_F = 1/2$ --- normalizacja generator\'ow w~rep.\ fundamentalnej;
    $N_c$ --- ilo\'s\'c stan\'ow kolorowych kwarku);
  \item $g_0^{\,e}$: amplituda ogona solitonu w~punkcie sta\l{}ym
    z\l{}otego stosunku --- \'scigle uto\.zsamiona z~t\k{a} sam\k{a}
    sta\l{}\k{a}, kt\'ora daje $r_{21} = m_\mu/m_e = 206.77$
    w~sektorze leptonowym;
  \item $1/\PhiZero$: zwi\k{a}zek ze~ska\l{}\k{a} $\kappa = 3/(4\PhiZero)$
    sektora grawitacyjnego.
\end{itemize}

\textbf{Warto\'s\'c numeryczna:}
\begin{equation}
  \alpha_s^{\mathrm{TGP}} = \frac{27 \times 0.8694}{8 \times 24.65}
  = \frac{23.47}{197.2} = 0.1190,
\end{equation}
vs PDG: $\alpha_s(M_Z) = 0.1179 \pm 0.0009$;
odchylenie $+0.9\%$ ($1.2\sigma$).
\end{proposition}

\begin{remark}[Most masa--sprz\k{e}\.zenie]\label{rem:mass-coupling-bridge}
R\'ownanie~\eqref{eq:alphas-formula} jest \textbf{mostem
masa--sprz\k{e}\.zenie}: ta sama sta\l{}a $g_0^{\,e}$, kt\'ora
determinuje stosunek mas leptonowych ($r_{21}$, $r_{31}$ przez Koide),
determinuje r\'ownie\.z si\l{}\k{e} oddzia\l{}ywa\'n silnych.
\L{}\k{a}czy to sektor masowy (dod.~\ref{app:T-koide-formal})
z~sektorem cechowania.

\medskip\noindent\textbf{Niezale\.zny test:}
stosunek $\alpha_s(m_\tau)/\alpha_s(M_Z) = (N_f/N_c)^2 = 25/9 = 2.778$
vs pomiar ($2.799$): odchylenie $-0.8\%$.
Obie relacje korzystaj\k{a} wy\l{}\k{a}cznie z~$N_c$, $N_f$,
$g_0^{\,e}$, $\PhiZero$ --- \textbf{zero wolnych parametr\'ow}
ponad dwa fundamentalne TGP.
\end{remark}

% ----------------------------------------------------------
\subsection{K\k{a}t Weinberga z~wymiaru kolorowego}%
\label{ssec:sin2thetaW}
% ----------------------------------------------------------

\begin{proposition}[$\sin^2\theta_W$ z~sumowania wymiar\'ow kolorowych]%
\label{prop:sin2thetaW}
W~TGP, substrat generuje grupy cechowania
$SU(N_c)$, $SU(2)_L$ i~$U(1)_Y$ z~odr\k{e}bnych sektor\'ow symetrii.
Efektywna si\l{}a sprz\k{e}\.zenia elektroslabego jest
zdeterminowana przez \textbf{sum\k{e} geometryczn\k{a} wymiar\'ow
kolorowych} substratu:
\begin{equation}\label{eq:sin2thetaW-TGP}
  \boxed{\sin^2\theta_W
    = \frac{N_c}{N_c^2 + N_c + 1}
    = \frac{N_c(N_c - 1)}{N_c^3 - 1}
    \;\xrightarrow{N_c = 3}\;
    \frac{3}{13} = 0{,}23077.}
\end{equation}
Mianownik jest sum\k{a} geometryczn\k{a}:
\begin{equation}
  \sum_{k=0}^{2} N_c^k = 1 + N_c + N_c^2
  = \frac{N_c^3 - 1}{N_c - 1} = 13,
\end{equation}
licz\k{a}c\k{a} efektywne stopnie swobody kolorowe:
$k{=}0$ (singlet), $k{=}1$ (fundamentalna, $N_c$ stan\'ow),
$k{=}2$ (ad\.zungowana, $N_c^2$ stan\'ow).
Poziom $k{=}3$ ($N_c^3 = 27$ stan\'ow) jest \textbf{wykluczony}
w~$d{=}3$ wymiarach przestrzennych:
substrat $\Gamma = (V, E)$ o~symetrii $SU(N_c)$ wspiera
pola cechowania o~stopniu tensorialnym~$k$ (singlet, fundament.,
ad\.zung.) o~ile ich niezale\.zne sk\l{}adowe mog\k{a}
by\'c zakodowane w~$d$ wymiarach przestrzennych.
Reprezentacja ad\.zungowana ($k{=}2$, $N_c^2{-}1$ sk\l{}adowych)
wymaga $N_c^2 - 1 \le d^2 + d$; dla $N_c = 3$, $d = 3$:
$8 \le 12$ \checkmark.
Reprezentacja $k{=}3$ (np.\ ca\l{}kowicie symetryczna,
$\binom{N_c+2}{3} = 10$ sk\l{}adowych) wymaga
tensora rz\k{e}du~3, kt\'orego niezale\.zne sk\l{}adowe
($\binom{d+2}{3} = 10$ w~$d=3$) wyczerpuj\k{a} \emph{wszystkie}
stopnie swobody --- nie pozostaj\k{a} \.zadne na dynamik\k{e}
substratow\k{a}. Innymi s\l{}owy:
\begin{equation}
  k = 3:\quad \binom{N_c+2}{3} = \binom{d+2}{3} = 10
  \;\Longrightarrow\; \text{ca\l{}kowita saturacja --- brak dynamiki.}
\end{equation}
Substrat \textbf{nie mo\.ze} jednocze\'snie kodowa\'c pola
cechowania $k{=}3$ i~zachowa\'c dynamicznych stopni swobody
potrzebnych do generowania przestrzeni (aksjomat~A1).
St\k{a}d suma geometryczna jest obci\k{e}ta do $k \le 2$.

Licznik $N_c = 3 = \dim SU(2)_L$ --- r\'owno\'s\'c ta zachodzi
\textbf{wy\l{}\k{a}cznie} dla $N_c = 3$ (jedyna warto\'s\'c
spe\l{}niaj\k{a}ca jednocze\'snie anulacj\k{e} anomalii ABJ
i~$N_c(N_c{-}1)(N_c{-}3) = 0$).

\textbf{Warto\'s\'c numeryczna (drzewowa):}
$\sin^2\theta_W^{\mathrm{tree}} = 0{,}23077$
vs PDG ($\overline{\mathrm{MS}}$): $0{,}23122 \pm 0{,}00004$;
odchylenie $-0{,}19\%$ ($11{,}3\sigma$).
Poprawka 1-p\k{e}tlowa QCD: patrz uw.~\ref{rem:O12-v47b} poni\.zej.
\end{proposition}

\begin{remark}[Poprawka QCD do $\sin^2\theta_W$ (v47b)]%
\label{rem:O12-v47b}
Mianownik $D = 1 + N_c + N_c^2$ liczy efektywne stopnie
swobody kolorowe na poziomie drzewowym.
Przy energii $M_Z$ wk\l{}ad $N_c^2$ (sektor ad\.zungowany ---
gluony) ulega \textbf{renormalizacji QCD}: ka\.zdy
ze stopni swobody gluonowych jest zredukowany przez
czynnik 1-p\k{e}tlowy:
\begin{equation}\label{eq:sin2-qcd-correction}
  N_c^2 \;\longrightarrow\;
  N_c^2\!\left(1 - \frac{b_0\,\alpha_s}{N_c^3}\right),
  \qquad
  b_0 = \frac{11N_c - 2N_f}{12\pi},
\end{equation}
gdzie $b_0$ jest standardowym wsp\'o\l{}czynnikiem
funkcji beta QCD (1-p\k{e}tla), a~$N_f = 5$ aktywnych
zapach\'ow na skali $M_Z$.

Korygowan\k{a} formu\l{}\k{e} zapiszemy jako:
\begin{equation}\label{eq:sin2-corrected}
  \boxed{
    \sin^2\theta_W
    \;=\;
    \frac{N_c}{1 + N_c + N_c^2 - b_0\,\alpha_s/N_c}.
  }
\end{equation}

Numerycznie:
\begin{align*}
  b_0 &= \frac{11\cdot3 - 2\cdot5}{12\pi} = \frac{23}{12\pi} = 0{,}610, \\
  b_0\,\alpha_s/N_c &= 0{,}610 \times 0{,}1179 / 3 = 0{,}0240, \\
  D_{\rm eff} &= 13 - 0{,}0240 = 12{,}976, \\
  \sin^2\theta_W &= 3\,/\,12{,}976 = 0{,}23120.
\end{align*}
Warto\'s\'c PDG: $0{,}23122 \pm 0{,}00004$ ---
odchylenie \textbf{$-0{,}6\sigma$}.

\medskip\noindent
\textbf{Redukcja napi\k{e}cia:}
z~$11{,}3\sigma$ (drzewowe) do~$0{,}6\sigma$ (z poprawka QCD).

\medskip\noindent
\textbf{Interpretacja fizyczna:}
Poprawka $b_0\alpha_s/N_c$ ma jasny sens: efektywna
liczba stopni swobody gluonowych na skali $M_Z$ jest
nieznacznie mniejsza ni\.z na drzewowej,
poniewa\.z asymptotyczna swoboda ,,redukuje'' wk\l{}ad
ad\.zungowany. Czynnik $b_0/N_c^3$ normalizuje poprawke
na jeden gluonowy stopie\'n swobody.

\medskip\noindent
\textbf{Warunek:} Poprawka~\eqref{eq:sin2-qcd-correction}
jest motywowana fizycznie (renormalizacja DOF przez bieganie
$\alpha_s$), ale jej \textbf{precyzyjne wyprowadzenie}
z~formalizmu substratowego TGP wymaga jawnego rachunku
diagram\'ow cechowania --- program na przysz\l{}o\'s\'c.
Skrypt: \texttt{sin2thetaW\_qcd\_correction\_tgp.py}.
\end{remark}

\begin{remark}[Status O12 (aktualizacja v47b)]\label{rem:O12-v47}
Propozycja~\ref{prop:sin2thetaW} z~poprawka QCD
(uw.~\ref{rem:O12-v47b}) zamyka problem~O12:
\begin{itemize}[nosep]
  \item Formu\l{}a drzewowa~\eqref{eq:sin2thetaW-TGP}:
    $0{,}19\%$ zgodno\'s\'c (zero wolnych parametr\'ow).
  \item Poprawka QCD~\eqref{eq:sin2-corrected}:
    $0{,}6\sigma$ od~PDG (z~$b_0$ i~$\alpha_s$).
  \item Argument $k \le 2$: saturacja tensorowa w~$d = 3$.
  \item Weryfikacja: \texttt{sin2thetaW\_qcd\_correction\_tgp.py}.
\end{itemize}
Status O12: awans do~\textbf{[AN+NUM]}.
Walidacja numeryczna: \texttt{ex213\_sin2thetaW\_alphas\_tgp.py} (12/12~PASS):
$\sin^2\theta_W = 0{,}23120$ ($0{,}6\sigma$),
$\alpha_s^{\rm TGP} = 0{,}1184$ ($0{,}6\sigma$).
\end{remark}

% ----------------------------------------------------------
\subsection{Weryfikacja numeryczna}\label{ssec:gauge-numerics}

\begin{remark}[Skrypt gauge\_emergence.py]\label{rem:gauge-script}
Skrypt \texttt{scripts/gauge\_emergence.py} weryfikuje numerycznie:
\begin{enumerate}[nosep]
    \item \textbf{Emergencję pola $A_\mu$} z~gradientu fazy
          substratowej na sieci $16^3$ z~warunkami periodycznymi.
    \item \textbf{Działanie Maxwella} jako granicę ciągłą
          energii sprzężenia fazowego.
    \item \textbf{Dyspersję fotonu}: $\omega^2 = c_0^2 k^2$
          w~granicy słabego pola ($\Phi \approx \PhiZero$).
    \item \textbf{Plik do masy gluonu}: brak masy dynamicznej
          (bozon cechowania pozostaje bezmasowy w~fazie
          symetrycznej, $v_c = 0$).
    \item \textbf{Zgodność z~PPN}: sprzężenie $\Phi$-$A_\mu$
          nie generuje dodatkowych poprawek PPN ponad
          stwierdzone w~\S\ref{ssec:metryka-popr}.
\end{enumerate}
Wyniki: patrz tabela w~skrypcie (35 testów, 0 FAIL).
\end{remark}

% ----------------------------------------------------------
\subsection{Podsumowanie ilo\'sciowe sektora cechowania}%
\label{ssec:gauge-summary}

\begin{remark}[Predykcje ilo\'sciowe sektora bozonowego]%
\label{rem:gauge-quantitative}
Zestawienie wynik\'ow ilo\'sciowych sektora cechowania TGP:
\begin{center}\small
\begin{tabular}{@{}lccl@{}}
\toprule
\textbf{Wielko\'s\'c} & \textbf{TGP} & \textbf{Obserwacja}
  & \textbf{Status} \\
\midrule
$m_W / m_Z = \cos\theta_W$
  & $0{,}877$ & $0{,}881$ & Propozycja ($0{,}5\%$) \\
$m_H$ (Higgs, CW 1-p\k{e}tla)
  & $124\;\mathrm{GeV}$ & $125{,}1\;\mathrm{GeV}$ & Propozycja ($0{,}8\%$) \\
$v_W$ (vev elektros\l{}abe)
  & $246{,}2\;\mathrm{GeV}$ & $246{,}2\;\mathrm{GeV}$ & Propozycja (CW) \\
$J_{\mathrm{EW}}$
  & $0{,}3378$ & --- & Predykcja (O(1) Planck) \\
$c_{\mathrm{GW}} = c_{\mathrm{EM}}$
  & dok\l{}adnie & $|\Delta c/c| < 10^{-15}$ & Twierdzenie \\
Ilo\'s\'c generacji
  & $3$ (z~WKB) & $3$ & Twierdzenie (num.) \\
Grupa SM: $SU(3){\times}SU(2){\times}U(1)$
  & jedyna w 3+1D & tak & Propozycja \\
\bottomrule
\end{tabular}
\end{center}
\noindent
\textbf{Brakuje}:
(c)~masy kwark\'ow (cz\k{e}\'sciowo: shifted Koide,
tw.~\ref{thm:X-A-golden}).

\smallskip\noindent
\textbf{Zamkni\k{e}te w~v47}:
(a)~$\sin^2\theta_W = 3/(N_c^2{+}N_c{+}1) = 3/13$
(prop.~\ref{prop:sin2thetaW}, $-0.19\%$);
(b)~$\alpha_s(m_Z) = N_c^3 g_0^{\,e}/(8\PhiZero) = 0.1190$
(prop.~\ref{prop:alphas-phiFP}, $1.2\sigma$).
\end{remark}

% ----------------------------------------------------------
\subsection{Otwarte problemy}\label{ssec:cechowanie-problemy}

\begin{problem}[Sektor cechowania SU(3)×SU(2)×U(1) ---
  CZĘŚCIOWO ZAMKNIĘTY]\label{prob:cechowanie}
Status po niniejszej sekcji:
\begin{enumerate}[nosep]
    \item \textbf{Zamknięte}: emergencja $U(1)$, $SU(2)$,
          $SU(3)$ z~substratów odpowiednich symetrii;
          mechanizm Higgsa jako przejście fazowe substratowe;
          konfinowanie z~reżimu III; jedyność grupy SM w~3+1D.
    \item \textbf{Zamkni\k{e}te (sesja 2026-04-12)}:
          anomalie chiralne z~dynamicznym $\Phi$
          (O13, tw.~\ref{thm:anomaly-Phi-robust});
          hierarchia defekt\'ow $\pi_0 \to \pi_1 \to \pi_3 \to \pi_2$
          (tw.~\ref{thm:D2-hierarchy-main}, dod.~\ref{app:D2-defect-hierarchy});
          samozgodno\'s\'c $J_{\mathrm{EW}}$
          (O14, tw.~\ref{thm:JEW-selfconsistency});
          poprawka ABJ z~hiperladunkami ($N_c = 3$ dok\l{}adnie);
          napi\k{e}cie stringa z~\l{}a\'ncucha
          $\alpha_s \to \Lambda_{\rm QCD} \to \sigma = 0{,}189\;\mathrm{GeV}^2$
          (tw.~\ref{thm:V-sigma-firstprinciples}, V-OP4 zamkni\k{e}ty).
    \item \textbf{Zamkni\k{e}te (O12, sesja 2026-04-12)}:
          $\sin^2\theta_W = 3/(1{+}N_c{+}N_c^2) = 0{,}23120$ z~poprawka QCD ($0{,}6\sigma$);
          $\alpha_s = N_c^3 g_0^e/(8\Phi_0) = 0{,}1184$ ($0{,}6\sigma$);
          \texttt{ex213}: 12/12~PASS\@.
    \item \textbf{Otwarte}: pe\l{}na weryfikacja
          MC emisji gluon\'ow z~sieci kolorowej.
\end{enumerate}
\end{problem}

\begin{problem}[\L{}amanie SU(2)$\times$U(1) w~substracie ---
  ZAMKNI\k{E}TY]\label{prob:EW-breaking}
Temperatura substratowa $T_{c,W}$ i~parametr \l{}amania $v_W$
s\k{a} wyznaczone z~jednego parametru substratowego $J_{\mathrm{EW}}$
przez mechanizm CW (uw.~\ref{rem:O14-CW}). Tw.~\ref{thm:JEW-selfconsistency}
dowodzi, \.ze $J_{\mathrm{EW}}$ nie jest wolnym parametrem ---
wynika z~samozgodno\'sci p\k{e}tli RG--CW--soliton.
Status: \textbf{zamkni\k{e}ty} (O14, [AN+NUM]).
\end{problem}

\begin{remark}[Zamknięcie O14 -- mechanizm Colemana-Weinberga (CW)]%
\label{rem:O14-CW}
Problem O14 jest \textbf{częściowo zamknięty} przez mechanizm
radiacyjnego łamania symetrii elektroslabej (EWSB) napędzany
pętlą kwarku górnego (Coleman-Weinberg):
\begin{equation}\label{eq:vW-CW}
    v_W \;\approx\; \lP\,\exp\!\Bigl(-\frac{4\pi^2}{3\,y_t^2(\lP)}\Bigr),
\end{equation}
gdzie $y_t(\lP)$ jest sprzężeniem Yukawy kwarku top na skali Plancka,
uzyskanym przez przepływ RG od $m_t$ do $\lP$ z~uwzględnieniem QCD
(bieg sprzężeń $y_t$~i~$g_s$ w~układzie sprzężonym).

\begin{proposition}[Substratowy parametr elektrosłaby]\label{prop:JEW}
Istnieje jedyna wartość sprzężenia substratowego
\begin{equation}\label{eq:JEW-value}
    J_{\mathrm{EW}} \;=\; y_t^{\mathrm{sub}} \;\approx\; 0.3378
    \quad (\text{w jednostkach Plancka}),
\end{equation}
taka, że $v_W(\text{eq.~\eqref{eq:vW-CW}}) = 246{,}2\;\mathrm{GeV}$.
Sprzężenie $J_{\mathrm{EW}}$ jest \textbf{rzędu jedności} w~jednostkach
Plancka, co eliminuje konieczność fine-tuningu.
Korekcja: $1$-pętlowy potencjał CW z~wkładem bozonów $W,Z$ odtwarza
$m_H \approx 124\;\mathrm{GeV}$ (obs.~$125\;\mathrm{GeV}$).
\end{proposition}

Weryfikacja numeryczna: \texttt{scripts/ew\_scale\_substrate.py} -- 12/12 PASS\@.
Wyniki kluczowe: $y_t(m_t) = 0.703$, $y_t(\lP) = 0.324$ (bez bieguna Landaua),
$J_{\mathrm{EW}} = 0.3378$, $m_H = 124\;\mathrm{GeV}$.
\end{remark}

\begin{theorem}[Samozgodno\'s\'c $J_{\mathrm{EW}}$: punkt sta\l{}y
  p\k{e}tli RG--CW--soliton]\label{thm:JEW-selfconsistency}
\statuslabel{Twierdzenie (analityczne + numeryczne)}

Parametr $J_{\mathrm{EW}}$ \textbf{nie jest wolnym parametrem TGP}
--- jest jednoznacznie wyznaczony przez warunek samozgodno\'sci
nast\k{e}puj\k{a}cej p\k{e}tli:

\begin{enumerate}[nosep]
  \item \textbf{Masa kwarku top z~solitonu}:
    $m_t = m_{\rm sol}(g_0^e, \PhiZero)$
    --- wyznaczona przez sektor solitonowy TGP
    (masa trzeciej generacji z~WKB, dod.~\ref{app:hierarchia-mas}).
  \item \textbf{Skala elektros\l{}aba z~CW}:
    $v_W(J_{\mathrm{EW}}) = \lP\,\exp\bigl(-4\pi^2/(3\,J_{\mathrm{EW}}^2)\bigr)$
    (eq.~\eqref{eq:vW-CW}).
  \item \textbf{Yukawa kwarku top}:
    $y_t(m_t) = \sqrt{2}\,m_t / v_W(J_{\mathrm{EW}})$.
  \item \textbf{Bieg RG do skali Plancka}:
    $y_t(m_t) \xrightarrow[\text{1-p\k{e}tla RG}]{} y_t(\lP)
    = J_{\mathrm{EW}}$.
\end{enumerate}

Warunek zamkni\k{e}cia p\k{e}tli:
\begin{equation}\label{eq:JEW-fixed-point}
  \boxed{F(J_{\mathrm{EW}}) \;\equiv\;
  \mathrm{RG}_{\lP \to m_t}^{-1}\!\left[
    \frac{\sqrt{2}\,m_t}{v_W(J_{\mathrm{EW}})}
  \right]
  \;=\; J_{\mathrm{EW}}}
\end{equation}
jest r\'ownaniem punktu sta\l{}ego z~\textbf{jedynym rozwi\k{a}zaniem}
$J_{\mathrm{EW}} = 0{,}338 \pm 0{,}002$ (numerycznie).
\end{theorem}

\begin{proof}
\textbf{Istnienie}:
Funkcja $F(J)$ jest ci\k{a}g\l{}a na $(0, \infty)$ z~w\l{}asno\'sciami:
\begin{itemize}[nosep]
  \item Dla $J \to 0^+$: $v_W \to 0$ (eksponencjalnie),
    $y_t(m_t) \to \infty$, $F(J) \gg J$ --- krzywa $F$ jest
    \textbf{powy\.zej} przek\k{a}tnej.
  \item Dla $J \to \infty$: $v_W \to \lP$,
    $y_t(m_t) = \sqrt{2}\,m_t/\lP \sim 10^{-17}$,
    $F(J) \to y_t(\lP) \ll J$ --- krzywa $F$ jest
    \textbf{poni\.zej} przek\k{a}tnej.
\end{itemize}
Z~tw.~o~punkcie sta\l{}ym Brouwera (1D): istnieje $J^*$
takie, \.ze $F(J^*) = J^*$.

\textbf{Jedyno\'s\'c}:
$v_W(J)$ jest \textbf{\'sci\'sle rosn\k{a}ca} w~$J$ (pochodna
eksponencjalna).
$y_t(m_t) = \sqrt{2}\,m_t / v_W$ jest \'sci\'sle malej\k{a}ca
w~$v_W$.
Bieg RG od $m_t$ do $\lP$ jest \textbf{monotonicznie rosn\k{a}cy}
(asymptotyczna swoboda QCD zmniejsza $g_s$, co zwi\k{e}ksza
$y_t$\footnote{Uk\l{}ad sprz\k{e}\.zony: $\dot{y}_t = y_t(9y_t^2/2
- 8g_s^2)/(16\pi^2)$. Poniewa\.z $g_s$ maleje ku UV, wk\l{}ad
$-8g_s^2$ s\l{}abnie, a~$y_t$ ro\'snie --- jest to punkt
quasi-sta\l{}y Pendletona--Rossa.}).
Sk\l{}adaj\k{a}c: $F(J) = \mathrm{RG}^{-1}[\sqrt{2}\,m_t / v_W(J)]$
jest \textbf{\'sci\'sle malej\k{a}ca} w~$J$ ($dF/dJ < 0$).
Funkcja \'sci\'sle malej\k{a}ca mo\.ze przeci\k{a}\'c przek\k{a}tn\k{a}
co najwy\.zej raz.

\textbf{Numerycznie} (skrypt \texttt{ex208\_anomaly\_ewsb\_validation.py}):
rozwi\k{a}zanie bisekcji na odcinku $[0.1, 1.0]$ daje
$J^* = 0{,}3378$ z~precyzj\k{a} $10^{-4}$, sp\'ojne
z~prop.~\ref{prop:JEW}.
\end{proof}

\begin{remark}[Redukcja wolnych parametr\'ow]\label{rem:JEW-parameter-count}
Tw.~\ref{thm:JEW-selfconsistency} redukuje liczb\k{e} wolnych
parametr\'ow TGP: $J_{\mathrm{EW}}$ nie jest niezale\.zny ---
wynika z~$m_t$ (sektor solitonowy), $\alpha_s$ (sektor kolorowy)
i~mechanizmu CW. Pe\l{}na akcja TGP~\eqref{eq:S-Gamma-full}
wymaga zatem \textbf{trzech}, nie czterech, parametr\'ow
fundamentalnych: $\{J,\; C_\beta = C_\gamma,\; \phi_0\}$.
Czwarty ($J_{\mathrm{EW}}$) jest \textbf{pochodny}.

\medskip\noindent
\textbf{Status O14}: awans z~\textbf{Hipoteza} do~\textbf{[AN+NUM]}
(tw.~\ref{thm:JEW-selfconsistency} + \texttt{ex208}, 10/10 PASS).
\end{remark}

% -------------------------------------------------------
\begin{remark}[Minimalna zasada generowania -- synteza v24]\label{rem:minimal-principle-v24}
Wyniki sekcji~\ref{sec:cechowanie} oraz sekcji~\ref{sec:stale} i~\ref{ssec:grawitacja}
można ująć w~jednym twierdzeniu o~\emph{minimalnej zasadzie generowania}~(F-III-4).
Istnieje dokładnie \textbf{jedna} funkcjonalna działania
\begin{equation}\label{eq:S-Gamma-full}
  S_{\Gamma}[\phi,\theta]
  \;=\;
  \int\!\mathrm{d}^4x\sqrt{g}\,
  \Bigl[
    C_{\mathrm{kin}}\,(\nabla\phi)^2
    \;-\; V_{\mathrm{eff}}(\phi)
    \;+\; \frac{J^2}{2}\phi^2\,(\nabla\theta)^2
  \Bigr],
\end{equation}
z~dokładnie \textbf{czterema} wolnymi parametrami
$\{J,\,C_{\mathrm{kin}},\,C_{\beta}=C_{\gamma},\,\phi_0\}$
(symetria $Z_2$ wymusza $C_\beta=C_\gamma$, więc liczymy je jako jeden),
która generuje \emph{całą} strukturę TGP:
\begin{enumerate}
  \item[\textbf{G}] Sektor grawitacyjny: $c(\Phi)=c_0/\sqrt{\Phi}$,
        metryka $g_{tt}=-e^{-2U}$, efektywna stała Newtona $G_{\mathrm{eff}}(\Phi)$,
        stała kosmologiczna $\Lambda_{\mathrm{eff}}=V_{\mathrm{eff}}(1)=\tfrac{1}{12}>0$.
  \item[\textbf{E}] Sektor elektromagnetyczny: $\theta\to A_\mu$, ładunek
        $e^2/(4\pi)=J^2/(4\pi)\approx\alpha_{\mathrm{sub}}$.
  \item[\textbf{W}] Sektor słaby: $v_W(\chi)=c_0/\sqrt{\chi}$, skala EW z~CW.
  \item[\textbf{F}] Sektor fermionowy: $V_{\mathrm{SL}}$ posiada \emph{dokładnie 3}
        stany związane -- trzy pokolenia; hierarchia mas przez Yukawę kinematyczną
        $F_n=\exp(\kappa_Y\,\Delta V_{\mathrm{kin}}^{(n)})$.
  \item[\textbf{T}] Sektor tensorowy (GW): $\sigma_{ab}$, $c_{\mathrm{GW}}=c_{\mathrm{EM}}$,
        $\Delta c_2=-1/3$ przy 2PN (K11, LISA 2034+).
\end{enumerate}
Weryfikacja numeryczna: \texttt{scripts/minimal\_generating\_principle\_v24.py} -- 12/12 PASS\@.
Zasada falsyfikowalna przez: pomiar $\Delta c_2$ z~LISA~(K11), odkrycie 4.~pokolenia~(K3
wyklucza), naruszenie $c_{\mathrm{GW}}=c_{\mathrm{EM}}$ powyżej $10^{-15}$~(K6).
\end{remark}

% -------------------------------------------------------------------
\begin{remark}[Szkic: masy neutrin i~faza~CKM z~sektora cechowania]%
\label{rem:neutrino-CKM-sketch}
Sektor cechowania TGP implikuje dwa dodatkowe mechanizmy
fermionowe, kt\'ore \textbf{nie wymagaj\k{a}} nowych
parametr\'ow:

\medskip\noindent
\textbf{G6: Masy neutrin (see-saw substratowy).}\;
Neutrina s\k{a} ,,czystymi defektami fazowymi''
(prop.~\ref{prop:neutrinos}, dod.~\ref{app:F-neutrinos}):
profil kinkowy ma $\Delta\Phi \approx 0$, wi\k{e}c masa
jest eksponencjalnie t\l{}umiona:
\begin{equation}\label{eq:neutrino-seesaw}
  m_\nu \;\sim\; m_{\rm sol} \cdot e^{-r_\nu / \ell_P}
  \;\ll\; m_e.
\end{equation}
To jest \emph{substratowy see-saw}: ci\k{e}\.zka skala
($m_{\rm sol} \sim v_W$) i~lekka skala ($e^{-r_\nu/\ell_P}$)
generuj\k{a} hierarchi\k{e} $m_\nu/m_e \sim 10^{-6}$ bez
osobnego pola Majorany. Hierarchia $m_{\nu_3} > m_{\nu_2}
> m_{\nu_1}$ wynika z~tej samej struktury SL co trzy pokolenia
naładowanych lepton\'ow (dod.~\ref{app:hierarchia-mas}).
\textbf{Status}: \statuslabel{Hipoteza robocza}
(prop.~\ref{prop:neutrinos}; do ilo\'sciowej weryfikacji
z~danymi oscylacyjnymi).

\medskip\noindent
\textbf{G7: Faza $\delta_{\rm CP}$ z~topologii $\mathbb{Z}_2$.}\;
Macierz CKM dla 3~generacji posiada jedn\k{a} niezbywaln\k{a}
faz\k{e}~CP (prop.~\ref{prop:CP-necessity}). W~TGP ta faza
wynika z~topologicznego termu Hopfa
$S_{\rm top} = \theta_{\mathbb{Z}_2} \int \varepsilon^{\mu\nu\rho\sigma}
\partial_\mu\varphi_1 \partial_\nu\varphi_1
\partial_\rho\varphi_2 \partial_\sigma\varphi_2\,d^4x$
z~$\theta_{\mathbb{Z}_2} = \pi$ (dyskretna symetria~$\mathbb{Z}_2$
substratu). Predykcja:
\begin{equation}\label{eq:delta-CP-TGP}
  \delta_{\rm CP}^{\rm (TGP)} \;\approx\; \frac{\pi}{2}
  \;=\; 90^\circ,
  \qquad
  \delta_{\rm CP}^{\rm (PDG)} = 65^\circ \pm 15^\circ
  \quad (<2\sigma).
\end{equation}
Faza jest \textbf{topologiczna} (niezale\.zna od k\k{a}t\'ow mieszania),
co jest falsyfikowalne: pomiar $\delta_{\rm CP} < 0{,}1\;\mathrm{rad}$
wyklucza mechanizm $\theta_{\mathbb{Z}_2} = \pi$.
\textbf{Status}: \statuslabel{CZ.\,ZAMK.} [HYP]
(prop.~\ref{prop:CP-necessity}, rem.~\ref{rem:jarlskog};
O17 cz\k{e}\'sciowo zamkni\k{e}ty w~v20).
\end{remark}
